\chapter[Convergence and noise]{Convergence, noise and limits of learning}\label{ch:cIII3}
\chapterbrief
{$\Omega\subset\R^n$ open connected, $p_0\in\Omega$, $u$ real-analytic on $\Omega$ (hypothesis (H-an)); the datum at step $K$ is the jet $j_Ku=(\partial^\gamma u(p_0))_{|\gamma|\le K}$ (model DM2), or $u$ known up to a noise of norm $\delta$ (model DM-$\delta$). $R_{\le m}$ is the space of operators of degree $\le m$.}
{to prove what can and cannot be learned from a growing amount of data: identification in the limit holds with an explicit rule, but it never comes with a certificate, no learner can decide whether a law exists, and from noisy data one gets stability only relative to a finite level of the dictionary.}
{Dim.~A41--A43. They prove Chapter~\ref{PAP-ch:convergence} of the paper.}

\begin{notebox}{Notation for the whole chapter}
For $m\le K$ put
\[
N_{m,K}=\{Q\in R_{\le m}:\ \partial^\beta(Q(\partial)u)(p_0)=0\ \text{ for all }|\beta|\le K-m\}.
\]
It is computable from $j_Ku$: in matrix form, the kernel of the matrix $M_{m,K}(j)$ with entries $j_{\alpha+\beta}$ (rows $\beta$, columns $\alpha$). $\nu_m(u)=\max\{\mathrm{ord}_{p_0}(Qu):Q\in R_{\le m}\setminus\Ann(u)\}$.
\end{notebox}

\section{Identification in the limit}

\begin{dimcard}{A41}{CONV-T1}{Stabilisation by level (CONV-L1), and identification in the limit (CONV-T1)}{\proved}{thm:convt1}
Assume (H-an).
\begin{enumerate}[label=(\alph*)]
\item (CONV-L1) $N_{m,K+1}\subseteq N_{m,K}$; $\bigcap_{K\ge m}N_{m,K}=\Ann(u)\cap R_{\le m}$; and $\nu_m(u)<\infty$ with $N_{m,K}=\Ann(u)\cap R_{\le m}$ iff $K\ge K^*(m)=m+\nu_m(u)$.
\item (CONV-T1) The \emph{minimal-degree rule} $m^*(K)=\min\{m\le K:N_{m,K}\ne0\}$, $O_K=\gcd(N_{m^*(K),K})$ satisfies: if $\Ann(u)\ne0$ with minimal degree $d$, then $m^*(K)=d$ and $O_K=g_d=\gcd(\Ann(u)\cap R_{\le d})$ for $K\ge K^*(d)$; if $\Ann(u)=(P)$ then $O_K=P$; if $\Ann(u)=0$ then $m^*(K)\to\infty$.
\item With an a priori bound $B$ such that $\Ann(u)$ is generated in degree $\le B$, $\gcd(N_{B,K})=G_u$ for $K\ge K^*(B)$.
\item The output sequence $(O_K)$ takes at most $(d+1)(d+2)/2$ distinct values, for every $u$.
\item The time of convergence is not bounded by any function of $d$ (for $n=2$: $u=1+x^{N+2}e^{x^2}$ has $\Ann(u)=(\tau)$ and $K^*=N+2$).
\item (No correct certificate) For $n\ge2$ there is no $K$ and no $\{0,1\}$-valued function $c$ of $j_Ku$ with: $c(j_Ku)=1\Rightarrow\Ann(u)=(O_K)$ for every analytic $u$, and $c=1$ for some $u$ with non-trivial principal $\Ann$.
\end{enumerate}
\end{dimcard}

\begin{proofidea}
The conditions defining $N_{m,K}$ shrink as $K$ grows, and their limit is the annihilator because an analytic function with zero Taylor series vanishes. Finite dimensionality of $R_{\le m}$ makes the stabilisation happen at a finite step; but nothing says when, and any finite jet is shared with a polynomial, which kills certificates.
\end{proofidea}

\begin{proof}
\step{1}{(a) A decreasing chain that stabilises}
The conditions of $N_{m,K+1}$ contain those of $N_{m,K}$. $Q$ lies in all of them iff the Taylor series of $Qu$ at $p_0$ vanishes iff $Qu\equiv0$ (analyticity and connectedness). For finiteness put $W_m=R_{\le m}/(\Ann(u)\cap R_{\le m})$ and $F_k=\{[Q]:\mathrm{ord}_{p_0}(Qu)\ge k\}$: a decreasing chain of subspaces of a finite-dimensional space with intersection $0$, so $F_{k_0}=0$ for some $k_0$ and $\nu_m\le k_0-1$. Finally, for $Q\notin\Ann(u)$, $Q\in N_{m,K}\iff\mathrm{ord}(Qu)\ge K-m+1$.

\step{2}{(b) The rule}
For $m<d$, $\Ann(u)\cap R_{\le m}=0$ so $N_{m,K}=0$ for $K\ge K^*(m)$; for $m=d$ and $K\ge K^*(d)$, $N_{d,K}=\Ann(u)\cap R_{\le d}\ne0$. So $m^*(K)=d$ and $O_K=g_d$. If $\Ann(u)=(P)$, every non-zero element has degree $\ge\deg P$, so $\Ann\cap R_{\le d}=\C P$. If $\Ann(u)=0$, for fixed $m$ all $N_{m',K}$ with $m'\le m$ vanish for large $K$, so $m^*(K)>m$.

\step{3}{(c) and (d)}
(c) is (a) with $m=B$: $N_{B,K}=\Ann(u)\cap R_{\le B}$ generates $\Ann(u)$, and the gcd of a generating set of an ideal of a UFD is the gcd of the ideal. (d) For $K\ge1$, $N_{K,K}\ne0$ ($\dim R_{\le K}\ge3$, one linear equation), so $m^*(K)\le K$; $N_{m,K}=0\Rightarrow N_{m,K+1}=0$, so $m^*$ is non-decreasing; at fixed $m^*$, $N_{m,K+1}\subseteq N_{m,K}$ makes $\deg O_K$ non-decreasing, and equal degrees give associated $O_K$. So $O_K$ is a function of the pair $(m^*(K),\deg O_K)$, monotone in lexicographic order, with at most $(d+1)(d+2)/2$ values.

\step{4}{(e) Non-uniform time}
Let $f(x)=1+x^{N+2}e^{x^2}$. $\tau$ kills $u$. If $Q=\tau Q_1+q(\xi)\in\Ann(u)$ with $q\ne0$, then $q(\partial_x)f=0$ on an interval, so $f$ would be an exponential-polynomial (analyticity), contradicting the growth $e^{x^2}$; so $\Ann(u)=(\tau)$. Computing orders: $\nu_0=0$, and for $m=1$, $Qu=af+cf'$ with $f'=(N+2)x^{N+1}e^{x^2}+\cdots$ of order $N+1$, so $K^*(1)=N+2$.

\step{5}{(f) No certificate}
Suppose $c$ and $u$ with $\Ann(u)=(P)$, $P$ non-constant, $c(j_Ku)=1$. Let $u'=T_Ku$ be the Taylor polynomial of order $K$: $j_Ku'=j_Ku$ and $\Ann(u')$ contains all monomials of degree $K+1$, so $R_n/\Ann(u')$ is finite-dimensional, whereas $R_n/(P)$ is infinite-dimensional for $n\ge2$. So $\Ann(u')\ne(P)$, while the rule and $c$ see the same data: $c(j_Ku')=1$ would assert $\Ann(u')=(O_K(u'))=(P)$, false.

\whyworks{the data at each finite stage are shared by a polynomial, whose annihilator is a different ideal; hence no stage can be certified, even though from some stage on the rule is right.}
\end{proof}
\uses{A2, A7.}

\begin{example}{Convergence at an unpredictable time}
For $u=1+x^{N+2}e^{x^2}$ the rule outputs $O_K=1$ (``no law of degree $1$ found'') for $1\le K\le N+1$ and $O_K=\tau$ from $K=N+2$ on, for every $N$. Nothing in the first $N+1$ derivatives announces $\tau$.
\end{example}

\begin{dimcard}{A42}{CONV-T2}{No learner decides whether a law exists: analytic Gold theorem (CONV-T2, T2$'$)}{\proved}{thm:convt2}
Let $n\ge1$ and $(\varphi_K)_{K\ge0}$ be \emph{any} sequence of maps $\varphi_K:\R^{J_K}\to\{\text{law},\text{no law}\}$ on jets of order $\le K$ (no computability, continuity or measurability assumed). Then there is a real entire function $u$ on $\R^n$ such that $\varphi_K(j_Ku)$ does not converge to the correct answer (``law'' if $\Ann(u)\ne0$, ``no law'' if $\Ann(u)=0$). (T2$'$) For $n\ge2$ the same holds with $u$ restricted to the class $\{\Ann(u)\text{ principal non-zero}\}\cup\{\Ann(u)=0\}$; hence no learner identifies $\Ann$ on that class.
\end{dimcard}

\begin{proofidea}
Build $u$ as a limit of polynomials: at each stage, add Taylor coefficients, chosen outside a proper algebraic set, that make every operator of degree $\le k$ fail to annihilate $u$; but wait until the learner has said ``law'' on the current polynomial, which it must eventually do. The learner then says ``law'' infinitely often on a function with $\Ann=0$.
\end{proofidea}

\begin{proof}
\step{1}{Step zero}
If some polynomial $p$ has $\varphi_K(j_Kp)\not\to$ ``law'', take $u=p$ ($\partial_t^{\deg p+1}\in\Ann(p)$). Assume $\varphi_K(j_Kp)\to$ ``law'' for every polynomial $p$.

\step{2}{An avoidance lemma}
Let $g_0(z)=e^{|z|^2}$. (i) $Q(\partial)(h\,g_0)\ne0$ for every non-zero polynomial $h$ and $Q\ne0$: its top-degree part is $Q_{\rm top}(2z)h_{\rm top}\,e^{|z|^2}\ne0$. (ii) For $m,L$ put $N^{>L}_{m,K}(g_0)=\{Q\in R_{\le m}:\partial^\beta(Qg_0)(0)=0\ \text{for }L<|\beta|\le K-m\}$; the intersection over $K$ consists of the $Q$ with $Qg_0$ a polynomial of degree $\le L$, so $\tau^{L+1}Q\in\Ann(g_0)=0$ and $Q=0$. By finite dimension there is $\bar K(m,L)$ with $N^{>L}_{m,\bar K}(g_0)=0$: the submatrix of rows $|\beta|>L$ has full column rank, so some minor $\Delta\ne0$ at the jet of $g_0$. (iii) $\Delta$ is a real polynomial in the coefficients $y=(j_\gamma)_{L<|\gamma|\le\bar K}$, not identically zero, so its zero set has empty interior and there is $y^*$ with $|y^*_\gamma|\le\gamma!/|\gamma|!$ and $\Delta(y^*)\ne0$. Any analytic $u'$ whose jet has coefficients $y^*$ in orders $(L,\bar K]$ has $\Ann(u')\cap R_{\le m}=0$.

\step{3}{Construction and limit}
Build polynomials $p_k$ and integers $H_0<H_1<\cdots$ with $\varphi_{H_k}(j_{H_k}p_k)=$ ``law''. At step $k$ apply the lemma with $m=k$, $L=H_k$, add the block $y^*_k$ in orders $(H_k,L'_k]$, and by Step~1 pick $H_{k+1}>L'_k$ with ``law'' on the new polynomial. Put $u=\lim p_k$. Its Taylor coefficients satisfy $|c_\gamma|\le1/|\gamma|!$, so $u$ is entire and real. For each $k$, $j_{H_{k+1}}u=j_{H_{k+1}}p_{k+1}$, hence $\varphi_{H_{k+1}}(j_{H_{k+1}}u)=$ ``law'' infinitely often, while the block $y^*_k$ gives $\Ann(u)\cap R_{\le k}=0$ for every $k$: $\Ann(u)=0$. So $\varphi$ is wrong on $u$.

\step{4}{(T2$'$) The principal class}
Test functions $v=q+\varepsilon t^{N-1}e^{|x|^2}$ with $q$ a real polynomial of degree $\le M$, $N=M+2$, have $\Ann(v)=(\tau^N)$ (divide $Q$ by $\tau^N$; the remainder gives a polynomial plus $e^{|x|^2}$ times a polynomial, impossible unless each coefficient vanishes) and $j_Mv=j_Mq$. Redo the construction with $v_k=q_k+\varepsilon_kt^{N_k-1}e^{|x|^2}$, $\varepsilon_k=1/(N_k-1)!$, $N_k=L'_k+2$, which lies in the class and on which the learner must say ``law''; the tails occupy orders $>H_{k+1}$, so the limit $u$ is entire (the tails are dominated by $e^{|t|}e^{\sum|x_i|^2}$), the learner says ``law'' infinitely often, and $\Ann(u)=0$.

\whyworks{any rule that is eventually right on all polynomials can be fooled by a function that agrees with polynomials longer and longer yet is differentially transcendental.}
\end{proof}
\uses{A7, A41.}

\begin{caution}[title=Reading]
``$\Ann(u)\ne0$'' is \emph{verifiable in the limit} (the rule ``say law at step $K$ if $m^*(K)=m^*(K-1)$'' is eventually right when $\Ann\ne0$ and wrong infinitely often when $\Ann=0$), but not \emph{decidable in the limit}.
\end{caution}

\section{Noisy data}

\begin{dimcard}{A43}{CONV-T4}{Discontinuity (T3), conditional stability (T4), and a discrepancy rule (T5)}{\proved}{thm:convt45}
\textbf{(T3)} Let $X$ be a topological vector space of functions on $\Omega$ continuously embedded in $\mathcal D'(\Omega)$, in which polynomials are dense and $e^{|z|^2}\in X$ (for example $H^m(\Omega)$ on a bounded Lipschitz domain). Then $\{\Ann=0\}$ and $\{\Ann\ne0\}$ are both dense, and $u\mapsto\Ann(u)$ is discontinuous everywhere.

\textbf{(T4)} Let $\Omega$ be bounded Lipschitz, $\Phi(u):\C^M\to L^2$, $\lambda\mapsto\sum_{|\alpha|\le n_0}\lambda_\alpha\partial^\alpha u$, with $\Phi(u)\lambda^*=0$, $|\lambda^*|=1$ and second singular value $\sigma_2(\Phi(u))\ge\rho>0$.
\begin{enumerate}[label=(\alph*)]
\item If $\|\tilde u-u\|_{H^{n_0}}\le\delta$ and $\tilde\lambda$ minimises $\|\Phi(\tilde u)\lambda\|$ on $|\lambda|=1$, then $\sin\angle(\tilde\lambda,\lambda^*)\le2\delta/\rho$; the order $\delta/\rho$ is optimal.
\item If $\|\tilde u_0-u\|_{L^2}\le\delta$ and $\|u\|_{H^{n_0+s}}\le\mathfrak M$, the same holds with $\delta$ replaced by $C(2\delta)^{s/(n_0+s)}(2\mathfrak M)^{n_0/(n_0+s)}$.
\item (Misfit certificate) If $\|\tilde u-u\|_{H^{n_0}}\le\delta$ and $\sigma_1(\Phi(\tilde u))>\delta$, then $\Ann(u)\cap\mathcal D_{n_0}=0$.
\end{enumerate}

\textbf{(T5)} With $d\ge1$ the minimal degree of $\Ann(u)\ne0$ and $s_m=\sigma_1(\Phi_m(u))>0$ for $m<d$, the discrepancy rule $m(\delta)=\min\{m:\sigma_1(\Phi_m(\tilde u))\le\delta_m\}$ gives $m(\delta)=d$ and $\sin\angle(\tilde\lambda(\delta),P)\to0$ at rate $\delta$ (or $\delta^{s/(d+s)}$) when $\delta<\delta_0=\tfrac12\min_{m<d}s_m$; if $\Ann(u)=0$, $m(\delta)\to\infty$. The threshold $\delta_0$ cannot be bounded below from the data.
\end{dimcard}

\begin{proofidea}
Singular-value perturbation (Weyl) gives (a) and (c); interpolation between $L^2$ and $H^{n_0+s}$ via Stein's extension gives (b); the discrepancy rule picks the first level at which the smallest singular value is below the noise level; the threshold fails to be computable because a tiny smooth perturbation creates a law of low degree with tiny margin.
\end{proofidea}

\begin{proof}
\step{1}{(T3) Both sets are dense}
Polynomials are dense and have $\Ann\ne0$. For $\{\Ann=0\}$: the map $Q\mapsto Q(\partial)e^{|z|^2}|_\Omega$ from $R_{\le m}$ to $\mathcal D'(\Omega)$ is injective, so choose test functions $\psi_i\in C_c^\infty$ with $\det(\langle Q_\alpha(\partial)g_0,\psi_i\rangle)\ne0$. The map $v\mapsto\langle Q_\alpha(\partial)v,\psi_i\rangle$ is continuous on $X$, and $D(\varepsilon)=\det\langle Q_\alpha(\partial)(u+\varepsilon g_0),\psi_i\rangle$ is a polynomial in $\varepsilon$ not identically zero, so non-zero for all but finitely many $\varepsilon$. For each $m$ only finitely many $\varepsilon$ are bad; the union over $m$ is countable, so there are arbitrarily small $\varepsilon$ that are good for every $m$, and then $\Ann(u+\varepsilon g_0)=0$ (\emph{condensed}: the package \S2.1 writes the determinant argument in full).

\step{2}{(T4a) A singular-vector estimate}
Put $E=\Phi(\tilde u)-\Phi(u)$; by Cauchy--Schwarz $\|E\|\le\delta$. Write $\tilde\lambda=c\lambda^*+w$, $w\perp\lambda^*$, so $|w|=\sin\angle$. Then
\[
\rho|w|\le\|\Phi(u)w\|=\|\Phi(u)\tilde\lambda\|\le\|\Phi(\tilde u)\tilde\lambda\|+\delta\le\|\Phi(\tilde u)\lambda^*\|+\delta\le\|\Phi(u)\lambda^*\|+2\delta=2\delta .
\]
Optimality: for $A=\mathrm{diag}(0,\rho)$, $E=\begin{pmatrix}0&0\\\delta&0\end{pmatrix}$ the minimal singular vector of $A+E$ has sine $\delta/\sqrt{\rho^2+\delta^2}$.

\step{3}{(T4b) and (T4c)}
(b) $e=\tilde u-u$ has $\|e\|_{L^2}\le2\delta$, $\|e\|_{H^{n_0+s}}\le2\mathfrak M$. With Stein's universal extension operator and H\"older's inequality in Fourier variables with exponents $(n_0+s)/s$ and $(n_0+s)/n_0$,
\[
\int(1+|\xi|^2)^{n_0}|\hat f|^2\le\Big(\int|\hat f|^2\Big)^{\frac s{n_0+s}}\Big(\int(1+|\xi|^2)^{n_0+s}|\hat f|^2\Big)^{\frac{n_0}{n_0+s}},
\]
so $\|e\|_{H^{n_0}}\le C\|e\|_{L^2}^{s/(n_0+s)}\|e\|_{H^{n_0+s}}^{n_0/(n_0+s)}$; then (a). (c) Weyl: $|\sigma_1(\Phi(\tilde u))-\sigma_1(\Phi(u))|\le\|E\|\le\delta$, so $\sigma_1(\Phi(u))>0$.

\step{4}{(T5) The discrepancy rule}
In the branch $\|\tilde u-u\|_{H^d}\le\delta$: for $m<d$, $\sigma_1(\Phi_m(\tilde u))\ge s_m-\delta>\delta$ when $\delta<s_m/2$, so $m$ is not chosen; for $m=d$, $\sigma_1(\Phi_d(\tilde u))\le\|\Phi_d(\tilde u)q\|\le\delta$ for a unit $q\in\Ann\cap R_{\le d}$. So $m(\delta)=d$, and the angle estimate of (T4a) with the kernel of dimension $r$ gives $\mathrm{dist}(\tilde\lambda,\Ann\cap R_{\le d})\le2\delta/\rho_d$. If $\Ann(u)=0$, every $s_m>0$, so for each $M$ and small $\delta$ no level $m\le M$ is chosen.

\step{5}{The threshold is not computable from the data}
Let $p=t^dx_1^d\cdots x_{n-1}^d$; one checks $\Ann(p)\cap R_{\le d}=0$ (the lowest-degree part $Q_{\rm low}$ of $Q$ maps $p$ to a nonzero polynomial). For almost every $\varepsilon$ the same determinant argument as in Step~1 gives $\Ann(u+\varepsilon p)\cap R_{\le d}=0$: a small smooth multiple of $p$ destroys the law of degree $d$ and leaves a margin that tends to $0$ with $\varepsilon$. The package builds from this, for every $N$ and $\eta>0$, a function $u'$ with $u'-u\in C^\infty(\bar\Omega)$, $\|u'-u\|_{H^N}\le\eta$, $\Ann(u')\ne0$ and $\delta_0(u')\le\eta$ (\emph{condensed}: the explicit choice of $u'$ is in the package, \S2.3 item~3). Hence no function $\Psi(\tilde u,\delta)>0$ of the data alone is a lower bound for $\delta_0(v)$ over all admissible $v$.

\whyworks{noise can only be handled relative to a finite level of the dictionary and with a margin $\rho$; both the level and the margin are properties of the unknown $u$, not of the data.}
\end{proof}
\uses{A7, A41, A42.}

\begin{example}{The heat equation from noisy data}
For the heat kernel $u=(t+1)^{-1/2}e^{-x^2/(4(t+1))}$ (minimal degree $d=2$, $\Ann(u)=(\tau-\xi^2)$), the discrepancy rule returned $m(\delta)=2$ for $\delta=10^{-3}$ and $10^{-4}$ in a SymPy/numerical check, and the observed ratio $\sin\angle/(2\delta/\sigma_2)$ never exceeded $0.27$ in $4000$ random trials.
\end{example}
